À l'instant $t=0$, on place l'interrupteur $K$ en position 1 (Figure 1).\\
La figure 2 indique l'évolution temporelle de la tension $u_{C}$ aux bornes du condensateur, et celle de la tension $u_{R}$ aux bornes du conducteur ohmique.\\
Donnée : $R=32 k \Omega$\\
1.1. Parmi les courbes (1) et (2), quelle est celle qui représente $u_{C}(t)$ ? Justifier.\\
1.2. Déterminer graphiquement la valeur de la force électromotrice $E$.\\
1.3. La constante de temps $\tau$ est la durée au bout de laquelle la charge du condensateur atteint $63 \%$ de sa charge maximale.\\
Déterminer graphiquement la valeur de $\tau$.\\
1.4. Vérifier que $C=47 \mu F$.

\begin{figure}[h]
\begin{center}
  \includegraphics[alt={},max width=\textwidth]{f47096eb-5ae7-4c60-ab2f-f9822564c423-4_540_803_504_1192}
\captionsetup{labelformat=empty}
\caption{Figure 2}
\end{center}
\end{figure}

1.5. L'expression de la tension $u_{C}$ est $u_{C}(t)=E$. $\left(1-e^{-t / \tau}\right)$.

Recopier, sur votre copie, le numéro de la question et choisir la lettre correspondante à la proposition vraie.\\
L'expression numérique de la charge $q(t)$ (en coulomb) du condensateur est :

\begin{center}
\begin{tabular}{|c|c|c|c|}
\hline
$\mathbf{A}$ & $q(t)=6.10^{-4} \cdot\left(1-e^{-1,5 . t}\right)$ & $\mathbf{B}$ & $q(t)=2,82.10^{-4} \cdot\left(1-e^{-0,67 . t}\right)$ \\
\hline
$\mathbf{C}$ & $q(t)=2,82.10^{-4} \cdot\left(1-e^{-666,7 . t}\right)$ & $\mathbf{D}$ & $q(t)=6 .\left(1-e^{-0,67 . t}\right)$ \\
\hline
\end{tabular}
\end{center}

